\documentclass[../main.tex]{subfiles}
\begin{document}
%==========================================TIÊU ĐỀ TẬP========================================
\begin{table}[h]
  \begin{adjustwidth}{-1cm}{}
    \begin{tabular}{>{\centering\arraybackslash}p{6cm}|>{\raggedright\arraybackslash}p{12.5cm}}
      \multirow{5}{6cm}
      % Cột bên trái: ÉPISODE và số tập
      {\\[-40pt]
        \flaregothic\fontsize{20pt}{20pt}\selectfont \centering
        {\contour{errortwo}{\textcolor{black}{ÉPISODE}}}\color{black}\\
        \burbank\fontsize{72pt}{72pt}\selectfont
        \includegraphics[height=3cm]{../../../../Image/Book?/number/num3.png}
        \\[18pt]
        % \flaregothic\fontsize{14pt}{14pt}\selectfont
        % \color{epattr}BIENVENUE, \uline{\color{Banpire} Mel} !
        % \flaregothic\fontsize{18pt}{18pt}\selectfont
        % \color{epattr}ÉPISODE PERDU
        \color{black}
      }
       &
      % Dòng thứ nhất: tên tập tiếng Nhật
      {\fotskip\fontsize{18pt}{18pt}\selectfont \ruby{青}{あお}\ruby{上}{がみ}\ruby{高}{こう}\ruby{校}{こう}の\ruby{崩}{ほう}\ruby{壊}{かい}
      } \\[6pt]
      % Dòng thứ hai: tên tập tiếng Anh
       & {\cafeteria\fontsize{18pt}{18pt}\selectfont The Fall of Aogami High}                                            \\[4pt]
      % Dòng thứ ba: tên tập tiếng Việt
       & {\vnmsans\fontsize{18pt}{18pt}\selectfont Sự sụp đổ của Trường Trung học Aogami}                                \\[8pt]
      % Dòng thứ tư: Nhân vật xuất hiện
       & {\caviar{\fontsize{14pt}{14pt}\selectfont Track used: The Cathedral}}
      \\[6pt]
      % Dòng cuối cùng: Thời gian diễn ra của tập (thường sẽ là ngày phát hành)
       & {\continuum\fontsize{16pt}{16pt}\selectfont Ngày phát hành: 05/03/2022}                                         \\[8pt]
    \end{tabular}
  \end{adjustwidth}
\end{table}
%===========================================NỘI DUNG==========================================
\color{black}

\begin{center}
  \pov{Haragen Kyuen}
  \uline{Hai ngày sau.\\Lớp 1-C - Trường trung học Aogami.}
\end{center}

Đã hai ngày kể từ sau cú đấm vô tình của Suzune.

Một cú đấm khá nhẹ, nếu so với cách mà Saikyou hay làm khi thúc giục tôi làm việc gì đó - những cú đó thì đôi lúc đau đến độ muốn nhập viện. Và tôi cũng từng phải chịu đựng đủ loại cú đấm, cú đạp như vậy từ đám bắt nạt trong suốt những ngôi trường mà hai anh em chúng tôi đã chuyển qua, nên\dots\ đau thì cũng chẳng đáng gì.

Saikyou cũng không trách Suzune-san, vì dẫu sao cú đấm đó nhẹ hều, giống như cách người ta chào hỏi thân mật hoặc tạo không khí thoải mái khi mới quen. Bình thường, chẳng có gì to tát.

Nhưng ở lớp 1-C này, hoặc ít nhất, với các học sinh của lớp nói chung, và các học sinh ở Trường Trung học Aogami nói riêng\dots\ mọi thứ không ``bình thường'' như thế.

Tôi nhớ rõ lúc đó, Touka-san đứng ngay cạnh, ánh mắt như muốn cắt đôi khoảng không giữa cô ấy và tôi. Cô ấy không hẳn nổi giận ầm ĩ, nhưng cái cách cô nghiêng đầu, hỏi ngắn gọn một câu với Suzune-san\dots\ nó mang một cảm giác khác hẳn. Không đơn thuần là bảo vệ tôi, mà như thể đang cố tuân theo một luật bất thành văn nào đó của lớp học này.

Người ta vẫn kể rằng Trường Trung học Aogami có một ``lời nguyền'', rằng những học sinh mới chuyển vào phải được đối đãi thật tốt. Thế nhưng, những hành động quá thân mật, quá gần gũi như vậy đôi khi sẽ kéo theo rắc rối không ai muốn gọi tên. Tôi từng nghe vài câu chuyện lặt vặt: một cái khoác vai vô tình dẫn đến cả tuần người kia không tới lớp, một lần nắm tay khiến hai người bỗng dưng nhìn nhau như kẻ xa lạ, hay một câu nói đùa xã giao tưởng chừng vô hại nhưng lại làm cho đối phương dằn mặt. Đúng là chuyện nghe qua thấy buồn cười, nhưng ở đây, nó lại nặng nề hơn người ta tưởng.

Và giờ, cú đấm nhẹ của Suzune-san cũng đã được xếp chung vào thứ ``không nên làm'' đó. Có lẽ họ cho rằng nó làm tôi khó chịu, từ việc tôi nhăn mặt sau cú đấm kia.

Tôi không rõ tại sao, nhưng cái không khí lạ lùng bao quanh lớp 1-C này\dots\ nó khiến mọi cử chỉ nhỏ nhất cũng trở thành một quân cờ trong ván cờ kỳ dị mà chúng tôi chưa từng chơi.

\ct

Giờ nghỉ giữa giờ.

Bầu không khí trong lớp vẫn nặng trĩu như bị nén lại, thứ không khí ấy đã kéo dài suốt hai ngày kể từ cú đập lưng vô tình của Suzune-san. Chuyện vốn chỉ là một hành động nhỏ, nhưng sau những lời đồn mà Inari-san đã buột miệng kể hôm kia, mọi thứ dường như bị phủ thêm một lớp sương mỏng, thứ sương khiến từng lời nói, từng nụ cười trở nên dè dặt. Đám học sinh vốn đã thưa vắng vì ``cúm mùa'' (theo như Touka-san giải thích hôm trước) giờ càng thưa thớt hơn.

Tôi cúi xuống, lấy sách vở từ trong cặp, chuẩn bị cho tiết học tiếp theo. Đầu óc tôi vẫn chưa thật sự thoát khỏi những suy nghĩ vẩn vơ về cái ``lời nguyền'' kia, dù tôi luôn tự nhủ rằng đó chỉ là mấy câu chuyện ma để hù dọa học sinh mới.

Nhưng\dots\ nhỡ đâu nó lại là câu chuyện có thật thì sao?

Liệu rằng có phải tôi là một trong những nguyên nhân khiến cho mọi thứ ở lớp này rối tung rối mù như thế?

Tôi cũng không thể chắc chắn được. Chỉ mong sao đó chỉ là những lời truyền miệng\dots

\dots

Bất chợt, một giọng nói nhỏ nhẹ nhưng có chút run run vang lên từ phía trước bàn: ``C-Chào cậu, Kyuen-san\dots''

Ngẩng đầu lên, tôi bắt gặp Suzune-san đang đứng đó. Ánh mắt cô ấy hơi cụp xuống, nụ cười cố gắng tỏ ra thân thiện nhưng lại cứng ngắc. Bên cạnh, Touka-san khoanh tay, thỉnh thoảng liếc sang Suzune như thể đang theo dõi sát một học trò vừa bị gọi lên bảng trả bài, gương mặt không giấu được sự căng thẳng.

\img{e1-3a}

``Touka-san, Suzune-san. Có chuyện gì sao?'' tôi mở lời.

Suzune-san, người mà tôi vẫn nghĩ tới như một ``quả cầu năng lượng'', giờ đây đang cười gượng, miệng lắp bắp: ``D-Dạo này cậu\dots\ thế nào rồi\~{}?''

Giọng nói như lẫn vào bầu không khí u ám, kèm theo những giọt mồ hôi nhỏ lăn nơi thái dương.

``Mình ổn. Có việc gì à, Suzune-san?'' tôi nghiêng đầu, rồi quay sang người đứng cạnh. ``Cả cậu nữa, Touka-san. Sao thế?''

``À\dots\ à thì\dots''

Cô bạn tóc xanh hít một hơi, rồi thở dài ra, dường như đã mất kiên nhẫn. Cô hích nhẹ vào hông cô bạn bên cạnh như ra hiệu: ``Suzune-san, cậu nên nói thẳng vào trọng tâm đi.''

``T-Tớ biết chứ\dots\ nhưng mà\dots''

``Sao? Có chuyện gì thì cứ nói thẳng. Cậu ấp úng thế này mình không đoán nổi đâu,'' tôi cố gắng cười nhẹ để bớt căng thẳng.

Suzune-san hít sâu, lấy hết dũng khí: ``N-Nhớ cái lúc mà\dots\ mình đánh cậu để giục cậu tham dự lễ hội ở thị trấn hôm kia chứ?''

Tôi ngẩn ra vài giây rồi nhớ ngay: ``À, cái vụ đó hả? Mình nói rồi mà, cái đấm đó không đau đâu. Không cần chấp--''

Chưa kịp nói hết, cô ấy đã cúi rạp đầu xuống rồi ngẩng lên, chắp tay trước ngực, giọng gấp gáp: ``\textbf{C-Cho mình xin lỗi!} Mình\dots\ vô thức cao hứng quá nên đập vào lưng cậu mạnh như thế! Mình không cố ý đâu!''

\img{e1-3b}

Tôi bật cười chát: ``Thì mình bảo rồi, không có gì đâu mà! Mình không để tâm đến nó đâu! Không cần làm nghiêm trọng vậy đâu!''

``H-Hành động đó\dots\ chỉ là một cách để mình\dots\ tiếp cận gần gũi hơn về thể xác thôi\dots!'' Suzune-san càng nói càng đỏ mặt.

``Mình biết chứ! Thật đấy, mình không bận tâm đâu! Mình từng bị nhiều cú đấm, cú thúc còn đau hơn nhiều rồi.''

Touka-san dường như thấy lời tôi vẫn chưa đủ để gỡ rối, liền xen vào, giọng hơi gấp: ``Có lẽ cậu không biết, nhưng Suzune-san đang tập rèn kỹ năng kết bạn\dots\ một trong số đó là tương tác gần gũi hơn về thể xác. Đôi lúc cậu ấy\dots\ hơi quá tay. Nhưng hoàn toàn không có ý làm hại đâu.''

Suzune-san gật lia lịa: ``Đ-Đúng vậy! Mình muốn gắn kết gần hơn với cậu\dots\ cậu hiểu chứ? Mình thực sự không cố ý đánh mạnh đâu!''

Nhìn gương mặt rối bời của hai người, tôi chợt thấy lạ. Bình thường tôi sẽ dễ dàng cho qua, nhưng từ hôm đó, từ khi nghe Inari-san kể về cái ``lời nguyền'' của trường\dots\ tôi lại thấy mình không thể chỉ cười rồi bỏ qua được.

Dù vậy, tôi vẫn giơ hai tay ra, cười trừ: ``Thật sự đấy, mình không để tâm đâu. Mình biết cậu làm vậy chỉ để thân thiết hơn với mình. Cú đấm đó chẳng phiền gì cả\dots\ thậm chí, mình còn thấy nó thể hiện tình cảm mà cậu dành cho mình nữa.''

Cả Touka và Suzune thở phào rõ rệt, vai trĩu xuống như vừa trút được gánh nặng. Tôi nghĩ nếu không nói vậy, chắc họ còn đứng đó lúng túng mãi.

``Vậy\dots\ sau tất cả\dots\ cậu sẽ tha lỗi cho Suzune-san chứ?'' Touka-san hỏi, ánh mắt hơi nghiêm.

``Cậu\dots\ sẽ bỏ qua cho mình chứ?'' Suzune-san nhìn tôi, đôi mắt ánh lên tia lo lắng, trực trào nước mắt.

``Tất nhiên rồi!'' tôi bật cười. ``Bạn bè thì đôi lúc cũng có những chuyện như thế mà. Mình hiểu cảm giác khi muốn kết thân với bạn mới, nên mình hiểu hành động của cậu lúc đó. Mình thật sự không bận tâm đâu, thậm chí còn thích cách tiếp cận đó là đằng khác!''

Hai cô gái như vừa được giải phóng khỏi một áp lực vô hình, mà tôi cảm nhận được rõ cả nghĩa bóng lẫn nghĩa đen. Tôi hiểu vì sao họ làm vậy, và cách tốt nhất là gật đầu bỏ qua. Thật lòng mà nói, tôi không ngại va chạm kiểu đó, chỉ là cú đấm tới quá bất ngờ và\dots\ tính hướng nội của tôi khiến mọi người dễ nghĩ tôi khó chịu.

Kể ra cũng tội cho Suzune-san. Chỉ vì một phút cao hứng mà liên lụy cả bạn thân.

``T-Tạ ơn giời\dots\ Kyuen-san tha lỗi cho mình rồi!'' Suzune thốt gần như reo lên, ngửa mặt lên trời.

\img{e1-3c}

``Thấy chưa?'' Touka mỉm cười, ``Chỉ cần làm rõ mọi chuyện và xin lỗi tử tế là được rồi.''

``Bạn cứ nghĩ quá làm gì.'' tôi nhún vai. ``Mình dễ tính lắm, mấy chuyện này mình cho qua hết.''

``Nhưng\dots\ nhưng mà\dots'' Suzune chợt cúi mặt, ``Sau những gì Inari-san nói hôm đó\dots\ mình sợ lắm.''

``À\dots\ cái lời đồn ấy à? Về học sinh mới chuyển đến?''

``Đ-Đúng! Chính nó! Mình chưa từng nghe bao giờ\dots\ Ngộ nhỡ chúng ta không bao giờ gặp lại nhau nữa, chỉ vì mình, thì sao\dots?''

Tôi cố giữ giọng bình thản: ``Mình nghĩ cậu không cần để tâm đâu. Chỉ là tin đồn truyền miệng thôi.''

Nhưng sâu trong lòng, một thoáng lạnh chạy dọc sống lưng. Tôi tự nhủ mình không tin vào mấy chuyện vô căn cứ ấy\dots\ nhưng giấc mơ đêm qua, nơi tôi thấy sân trường chìm trong ánh hoàng hôn kỳ lạ và bóng những người bạn mờ dần\dots\ lại khiến tôi không thể gạt bỏ hoàn toàn.

Touka-san chen vào, dứt khoát: ``Cậu không cần lo đâu. Tôi từng nghe về nó rồi, nhưng nó chỉ là bịa đặt vô căn cứ thôi.''

Tôi gật đầu\dots\ nhưng trong lòng vẫn thấy mơ hồ.

Tôi liếc sang cô em gái của tôi, người đang cùng với Hanabi-san và Saki-san chứng kiến mọi chuyện với những nét mặt e ngại. Tôi ra dấu hiệu ``OK'', ám chỉ rằng mọi chuyện đã được giải quyết êm xuôi. Saikyou khẽ cau mày nhưng rồi cũng gật đầu, như thể đã tạm an tâm, ít nhất là ở thời điểm hiện tại.

Câu lời nguyền đó\dots\ vẫn khiến tôi day dứt mãi: ``Nếu bạn không đối xử tốt với một học sinh mới chuyển trường, những điều bất hạnh sẽ xảy ra.''

Tôi không muốn tin vào mấy thứ kiểu này, nhưng cảm giác kỳ lạ từ giấc mơ trước ngày đầu tiên đến Aogami vẫn chưa buông tha tôi. Trong mơ, hành lang dài bất tận, tiếng bước chân lạ, và những dãy bàn trống trải\dots\ tất cả hòa vào nhau thành một áp lực vô hình. Giờ đây, nhìn Suzune và Touka lo lắng đến mức này, tôi lại cảm thấy ranh giới giữa ``tin đồn'' và ``thực tại'' mỏng manh hơn tôi nghĩ.

\pov{Haragen Saikyou}

Tôi đứng ở bên cửa sổ, mắt khẽ liếc sang phía bên kia. Nii-chan đang nói chuyện với Touka-san và Suzune-san, không khí giữa ba người trông\dots\ khó tả. Cô bạn tóc xanh vẫn giữ nét mặt bình tĩnh nhưng tôi để ý ánh mắt cô ấy có gì đó như đang kiểm soát tình hình, còn cô bạn tóc hai cục thì cúi đầu rối rít, nói gì đó mà tôi đoán là xin lỗi. Cú đấm vào lưng Nii-chan hai ngày trước tưởng chừng chẳng có gì  lại thành chuyện lớn thế này.

Tôi khẽ thở dài, lắc đầu như để tự mình xua bớt cảm giác căng thẳng trong không khí.

``Đúng là mọi chuyện rối bời thật\dots'' tôi lẩm bẩm, vừa đủ để Saki-san và Hanabi-san bên cạnh nghe thấy.

Saki-san ngước nhìn về phía bàn của anh tôi, khẽ nhún vai: ``Cũng chẳng thể trách được họ. Nhưng kể ra cũng tội Touka-san, tự dưng bị cô bạn thuở nhỏ ấy lôi vào\dots''

Hanabi-san cau mày, xen vào: ``Tất cả là do cái tin đồn ngớ ngẩn kia sao?''

Tôi ngẩng lên, nhớ lại buổi hôm ấy ở hai ngày trước, khi Inari-san bất chợt thốt ra chuyện ``lời nguyền'' đó, giọng mình vô thức chậm lại: ``Cái mà Inari-san vô thức thốt ra hôm đó à?''

``Ừm.'' cô bạn Năng nổ gật đầu chắc nịch.

Cô bạn tóc hai bím nghiêng đầu, nhìn Hanabi-san: ``Trông cậu có vẻ không sợ nhỉ. Cậu không tin vào lời nguyền đó sao, Hanabi-san?''

``Dĩ nhiên là không,'' cô ấy trả lời không chút do dự, rồi quay sang tôi. ``Cậu thì sao, Saikyou-san?

Tôi khựng lại. Một phần trí óc đang dõi theo Nii-chan ở bên bàn học ấy, một phần khác lại chợt thoáng hiện ra giấc mơ kỳ lạ của đêm trước ngày đầu tiên đến trường\dots\ cái cảm giác bất an len lỏi vẫn chưa biến mất.

``Saikyou-san?''

Tiếng gọi của cô ấy đưa tôi về thực tại. Tôi ấp úng đáp: ``Ể? Ờ\dots\ ừm\dots\ Không.'' Nhưng,  giọng của tôi thiếu đi sự chắc chắn thường thấy.

Nếu dựa vào những gì mà hai anh em chúng tôi mơ thấy vào đêm hôm đó, những mẩu báo, những bẩu ghi chú mà tôi đã đọc trong giấc mơ kia, lời nói của Lớp trưởng có lẽ khá hợp lý phần nào.

Nhưng thực sự, tôi vẫn chưa thể chắc chắn vào bản thân mình được. Tôi thoáng mỉm cười gượng: ``Cơ mà, Inari hoàn toàn bị thuyết phục rồi\dots''

``Cậu ấy là tuýp người nhát gan mà,'' Saki-san nói, như thể để trấn an cô bạn đang run rẩy ở đằng kia. Từ khi vụ việc đó xảy ra, hầu như chúng tôi luôn thấy Lớp trưởng cứ ôm đầu lại, thỉnh thoảng lặp đi lặp lại những câu nói khiến chúng tôi vừa lo lắng vừa sợ hãi.

Quả thật, Inari-san bị ám ảnh nặng tới mức đó cơ đấy. Có lẽ cô ấy biết gì đó về cái gọi là lời nguyền kia, nhưng\dots\ liệu cô ấy sẵn sàng chia sẻ không?

Hanabi-san khoanh tay, giọng cương quyết: ``Chuyện này không thể tiếp diễn được. Chúng ta phải làm gì đó để khiến cậu ấy thấy đó chỉ là một lời nguyền vô căn cứ.''

Tôi thoáng nghiêng đầu: ``\dots Nhưng mà bằng cách nào? Đúng hơn là, cái lời nguyền kia bắt nguồn từ đâu và từ lúc nào vậy?''

Saki-san gõ nhẹ ngón tay lên cằm, suy nghĩ: ``Mình cũng không chắc nữa\dots\ Nếu liên quan đến ngôi trường này, mình nghĩ ở thư viện trường sẽ có cái gì đó để ta có thể tìm hiểu đấy.''

``Thì thế,'' Hanabi-san đáp lại, rồi thở dài chán chường. ``Tôi cũng nghĩ như cậu đấy, nhưng khi lên thư viện thì chẳng tìm thấy được gì cả.''

``Không có tài liệu gì sao?'' tôi hỏi lại.

Cô chủ bán hoa nheo mắt: ``Nghe có vẻ sai sai thế nào ấy\dots''

Điều đó rất kỳ lạ. Không biết vì sao, nhưng tôi cảm thấy có một thế lực nào đó đã xóa sổ những gì liên quan tới nó trong trí nhớ của mọi người.

Cô bạn Năng nổ khoát tay, quyết định một cách chớp nhoáng: ``Vậy chốc nữa khi nào về sẽ hỏi ông mình vậy. Ông mình cũng từng học ở ngôi trường này mà. Thậm chí mình có thể hỏi mọi người quanh thị trấn này. Chắc chắn phải có ai đó biết về lời đồn ác ý này chứ.''

Saki-san gật đầu: ``Kể ra thì cũng lạ. Nếu chuyện tương tự đã từng xảy ra trước đây mà không có ghi chép gì, thì tại sao lại không có gì ghi lại nhỉ?''

Cô bạn nhún vai: ``Đó. Nếu như vậy thì ta cũng có thể thuyết phục Inari rằng lời đồn kia chỉ là tin vịt được rồi.''

``Ý hay đấy,'' gương mặt cô sáng bừng. Vậy mình sẽ về hỏi mẹ. Bà ấy cũng từng theo học ngôi trường này, giống như ông cậu vậy. Có khi bà ấy sẽ biết được gì đó.

Hanabi-san mỉm cười nhẹ: ``Nếu được thế thì tốt quá.''

Một thoáng yên lặng lướt qua giữa ba người, trước khi Hanabi-san xua tay, như muốn đẩy bớt những suy nghĩ nặng nề ra khỏi đầu: ``Trời ạ, mấy chuyện kiểu này cứ vẩn vơ trên đầu rồi ong ong bên tai kiểu này ghét ghê. Tốt nhất là ta phải làm rõ chuyện này càng sớm càng tốt.''

Tôi nhìn ra sân trường, nơi nắng chiều hắt ánh vàng xuống những tán cây, rồi khẽ đáp: ``Ừ\dots\ Tôi cũng hy vọng ai đó sẽ làm sáng tỏ điều này.''

``Chắc chắn rồi,'' Saki-san cười. ``Kỳ nghỉ hè cũng sắp tới rồi đó. Hy vọng sau vụ này chúng ta có thể thoải mái hơn\dots''

Cuộc trò chuyện với Hanabi-san và Saki-san tạm kết thúc khi cả ba đồng ý sẽ tìm hiểu thêm. Nhưng ngay cả khi những câu chữ cuối cùng vừa rời khỏi môi, tâm trí tôi đã trôi dạt đi nơi khác.

Khi ấy, không gian hành lang chiều muộn bỗng trở nên mờ nhòe, giống như trong giấc mơ đêm trước ngày đầu tiên tôi và Nii-chan đến trường. Những gì viết trên vài tờ báo cũ úa vàng, mép giấy đã rách và cong lại, một vài mẩu mà phần sau bị xé mất\dots\ tất cả đều trùng hợp tới khó tin.

Cảm giác khi ấy\dots\ như có ai đang thì thầm ngay sau gáy, nhưng quay lại thì không thấy ai.

Trong giấc mơ ấy, dù cố gắng tìm kiếm, chúng tôi cũng không thể lần ra bất kỳ mẩu thông tin nào nói rõ nguồn gốc của ``lời nguyền''. Những mẩu báo chỉ nhắc lửng lơ về ``một lời đồn gây náo loạn học sinh'', về chuyện ngôi trường bị đổ sập vì lý do không rõ, và rồi dừng ở đó, như thể có bàn tay vô hình đã chặn lại phần còn lại của câu chuyện. Càng đi, tôi càng thấy rõ hơn một điều: ai đó đã cố ý khiến chúng biến mất.

Khi trở lại thực tại, tôi chợt nhận ra lời Saki-san ban nãy: ``Nếu chuyện tương tự từng xảy ra, sao lại không có gì ghi lại?'' Nghe giống một lời ám chỉ hơn là thắc mắc.

Hanabi-san khi nói về việc ``hỏi ông'' cũng tránh nhìn thẳng vào mắt tôi, và ngay cả Nii-chan\dots\ mấy hôm nay, mỗi khi chủ đề này được nhắc đến, tất cả mọi người đều lảng sang chuyện khác. Cảm giác như cả thị trấn này đang đồng lòng giữ kín điều gì đó. Và nếu bí mật đó có liên quan đến ``lời nguyền'' thì e rằng sự thật còn khủng khiếp hơn những gì Inari-san tưởng tượng.

Tôi cắn môi, đưa mắt nhìn ra sân trường. Ánh nắng chiều đã gần tắt, kéo bóng những tán cây dài hun hút trên mặt đất. Một làn gió nhẹ thoảng qua, nhưng lại khiến tôi rùng mình.

Bất chợt\dots

*XOẸT*

Một tia sáng xanh nhạt thoáng lướt qua tầm mắt. Chỉ trong khoảnh khắc, trong lớp học, tôi thấy một bóng người đứng im lặng cách đó không xa, còn những người còn lại trong lớp 1-C thì biến mất\dots

\img{e1-3d}

Cô gái ấy\dots\ mái tóc nâu dài chạm lưng, đeo một cài tóc màu đỏ lệch sang trái. Ánh mắt xanh biếc nhìn về phía tôi, lạnh lùng đến mức khiến nhịp tim khựng lại. Trên người cô là bộ đồng phục đen gọn gàng, cùng chiếc khăn quàng xanh nổi bật hệt như điểm nhấn duy nhất giữa sắc tối trùm xuống.

Tôi đã gặp cô ấy rồi\dots\ không phải ở ngoài đời, mà là trong chính giấc mơ ấy. Cảm giác như một mảnh ký ức bị kéo ngược trở lại.

``Saikyou-san?'' giọng Hanabi-san vang lên, kéo tôi khỏi khoảng trống rợn ngợp ấy. Tôi chớp mắt. Khoảnh khắc vừa rồi tan biến, cả lớp học quay trở về như cũ.

``K-Không có gì\dots'' tôi đáp vội, nhưng vẫn quay người tìm kiếm. Không còn ai ở đó.

Ánh mắt tôi lập tức hướng về Nii-chan, người vừa kết thúc cuộc nói chuyện với Touka-san và Suzune-san. Tôi tiến lại gần và hạ giọng: ``Nii-chan\dots\ em vừa thấy một cô gái\dots\ tóc nâu dài đến lưng, mắt xanh, đồng phục đen và khăn quàng xanh.''

Anh sững lại. Chỉ một cái nhìn thoáng qua thôi cũng đủ để tôi biết: anh cũng đã thấy. Sắc mặt Nii-chan khẽ đổi, như thể vừa gặp lại thứ gì đó mà mình không muốn chạm đến lần nữa.

Không ai khác, ngoài chính \emph{ả ta.}

Một luồng lạnh buốt chạy dọc sống lưng. Hai anh em chúng tôi không cần phải nói gì thêm, vì cùng hiểu rằng\dots\ cô gái đó không chỉ tồn tại trong giấc mơ.

Tôi vẫn còn chưa kịp hoàn hồn sau cái bóng của cô gái tóc nâu thì một âm thanh trầm đục bất chợt nổi lên, nghe như vọng từ sâu trong lòng đất.

*Ù Ù Ù\dots*

Âm thanh ấy ban đầu mơ hồ, như tiếng ầm ì của cỗ máy khổng lồ đang khởi động ở đâu đó dưới chân mình. Không khí trong lớp trở nên nặng nề, những khung cửa kính khẽ rung lên lách cách, còn những chiếc bàn, ghế bắt đầu dịch chuyển từng nhịp nhỏ nhưng rõ rệt.

``Tiếng gì vậy?'' Touka-san ngẩng lên, đôi mày khẽ cau lại.

Những hạt bụi từ trần nhà rơi xuống, lơ lửng trong luồng sáng chéo nghiêng của buổi chiều. Tiếng kẽo kẹt khe khẽ vọng ra từ dầm gỗ trên cao, như tiếng thở dài của một tòa nhà già nua đang chống chọi lại sức nặng vô hình.

Khoảnh khắc toàn bộ phòng lớp học ấy rung lên, tôi chợt sững người lại.

\dots Quen lắm.

Như một cánh cửa trong đầu bật mở, ký ức ùa về, chính là khoảnh khắc đó. Khoảnh khắc trong giấc mơ đêm ấy, khi Shino xuất hiện từ trong một lớp học nhìn chúng tôi, đôi mắt xanh sắc lẻm, và ngay sau đó là tiếng gầm rú của trận động đất. Hành lang dài vô tận rung lắc dữ dội, rồi sụp đổ thành biển gạch vụn.

Không chỉ một lần\dots\ mà nhiều lần. Giấc mơ đó đã bắt chúng tôi sống lại cảnh ấy, lặp đi lặp lại như một bản án.

*\textbf{Ù Ù\dots\ Ù Ù!!}*

Cơn chấn động lên mạnh hơn. Những chai nước đặt trên bàn nghiêng đổ, vệt nước loang ra mặt gỗ. Ghế tự trượt trên sàn theo nhiều hướng khác nhau.

``\textbf{Đ-Động đất kìa!}'' Saki-san kêu lên, giọng lạc đi vì hoảng. ``Mọi người, mau nấp xuống bàn đi!''

Cả lớp rối loạn. Tiếng chân, tiếng ghế va nhau chan chát, tiếng người gọi nhau lẫn vào tiếng rạn nứt rợn người từ đâu đó trên trần. Những ô cửa sổ vỡ tung thành nhiều mảnh do chấn động, những hạt bụi từ trên trần dần rơi xuống như suối.

Hanabi-san quỳ xuống bên cạnh Inari-san, hai tay ôm chặt vai bạn mình. ``I-Inari! Chúng ta sẽ ổn thôi!''

Inari co người lại, ôm đầu, đôi mắt mở to đầy kinh hoàng. ``T-Tất cả là vì chúng ta đã làm phật lòng Kyuen-san rồi đó!'' giọng cô run rẩy, gần như nghẹn lại. ``\textbf{Cứu mình với, Hanabi-chan!}''

Tôi liếc về Nii-chan, nhưng anh chưa kịp đáp thì Hanacô bạn thuở nhỏ của cô ấy đã lắc mạnh vai Inari-san. ``\textbf{Cậu giữ bình tĩnh lại đã nào! Đừng có hoảng hốt như vậy!}''

Nhưng Inari không nghe. Cô lẩm bẩm như người mê sảng: ``Đây là sự trừng phạt\dots\ chắc chắn là vậy\dots!''

Tôi vừa định chen vào thì từ phía trên đầu họ, một mảng trần bắt đầu rạn. Những vết nứt chạy ngoằn ngoèo như mạng nhện, từng hạt bụi vỡ rơi lả tả.

``\textbf{\dots! Inari-san! Hanabi-chan! Cẩn thận phía trên!}'' tôi hét lên, tay chỉ thẳng lên trên.

*RẦM! XOẢNG! XOẢNG!*

Tất cả diễn ra quá nhanh. Khối trần lớn đổ ập xuống, kèm theo bê tông, gỗ và gạch vụn, nuốt chửng hai người trong đám bụi mù. Âm thanh va đập chát chúa khiến màng nhĩ tôi rung lên, còn tim tôi như rơi thẳng xuống dạ dày.

\img{e1-3e}

``\textbf{Hanabi-san! Inari-san!!!}'' tôi hét lên, lao về phía đống đổ nát kia nhằm ứng cứ, nhưng lập tức bị Nii-chan giữ chặt.

``Đừng có dại, Saikyou! Nguy hiểm đấy!''

``Nhưng mà\dots\ nhưng mà\dots!''

``Không nhưng nhị gì cả! Mau nấp xuống bàn, nhanh lên! Lo cho thân mình trước đi đã!''

Nói đoạn, anh ấy kéo tôi dúi xuống gầm bàn, cả hai chúng tôi khom lưng, co gối lại cho vừa khoảng trống chật hẹp.

Tôi ngước nhìn ra ngoài. Saki-san, Touka-san và Suzune-san đang bám chặt lấy nhau, cố lê từng bước giữa mặt sàn rung chuyển. Từng mảng vữa rơi xuống ngay sát họ, như những nhát búa giáng thẳng vào lớp học.

\img{e1-3f}

Cô bạn tóc xanh túm lấy tay cô bạn thuở nhỏ của mình, gấp gáp trấn an: ``C-Chúng ta sẽ ổn thôi mà!''

``Nhưng\dots\ Inari-san và Hanabi-san đã\dots'' Suzune-san nói đứt quãng, nước mắt trào ra, giọng run rẩy như sắp vỡ vụn.

``Họ sẽ không sao, mình tin là vậy! Mau lên!''

Nhưng Suzune-san càng khóc to hơn: ``Tất cả là lỗi của mình! Lỗi của mình hết! Nếu mình không mạnh tay lúc đó thì--!''

``\textbf{Giờ không phải lúc để dằn vặt bản thân đâu!}'' tôi hét lớn, chỉ vào một gầm bàn gần đó. ``\textbf{Mau trốn xuống đó rồi tính tiếp!}''

Saki-san gật đầu lia lịa: ``S-Saikyou-san nói đúng đó! Chúng ta cũng phải trốn thôi--''

*RẮC!*

Âm thanh khô khốc vang lên từ phía trần trên họ. Không kịp hét lên hay cảnh báo, một khối bê tông khổng lồ cùng dầm thép gãy rơi xuống, nghiền nát cả ba trong tiếng sập ầm ầm.

``\textbf{Saki-san! Touka-san! Suzune-san!!}''

``\dots Khỉ thật chứ! Saikyou, níu chặt lấy anh!''

``Nhưng mà họ\dots''

``Em định tính để anh bị đống đất đá kia nuốt chửng luôn sao!?''

``\dots!''

Không còn cách nào khác, tôi ôm chầm lấy Nii-chan, bất lực nhìn đống đổ nát dần vùi lấp những người bạn học cùng lớp mà không thể làm gì được.

Bụi dày đặc đến mức tôi không còn nhìn thấy hình dạng con người nữa. Cổ họng tôi nghẹn cứng, tim đập loạn xạ, và tay tôi siết lấy cạnh bàn đến tê dại.

Trận động đất vẫn chưa dừng lại. Ngược lại, nó mạnh hơn. Mặt sàn rung bần bật, cả căn phòng nghiêng hẳn sang một bên như bị nhấc lên rồi thả xuống.

Tiếng gãy vỡ nối tiếp nhau, tường nứt ra thành những khe lớn, mái ngói trượt khỏi xà, rơi ào ào như mưa đá. Từng khoảng sáng ngoài cửa sổ bị che khuất bởi bụi và mảnh vỡ.

Chỉ trong chớp mắt, cả lớp 1-C chìm trong biển gạch vụn, bàn ghế bị chôn vùi, tiếng kêu cứu cũng bị nuốt trọn. Và đâu đó, dưới cơn địa chấn kinh hoàng này\dots\ tôi có cảm giác mình lại đang ở trong giấc mơ. Giấc mơ nơi không chỉ lớp học này, mà cả ngôi trường Aogami, sụp đổ thành tro bụi.

\ct

\dots

\dots

Tiếng rung chấn cuối cùng cũng dần lịm đi, để lại trong không khí một thứ tĩnh lặng nặng trĩu đến mức cả tiếng tim tôi đập cũng trở nên chói tai. Cơn bụi mù dày đặc vẫn treo lơ lửng trong lớp, mùi vữa vụn, gỗ cháy và sắt gỉ trộn lẫn, chát đắng nơi cuống họng. Tôi khẽ mở mắt, nhận ra mình vẫn đang co ro dưới gầm bàn, toàn thân cứng đờ.

Từ phía bên ngoài, tôi còn nghe đâu đó tiếng còi cấp cứu và cứu hỏa bên ngoài, cùng với đó là những tiếng người bên ngoài:

``Trường Trung học Aogami bị sụp đổ rồi!''

``Các đồng chí mau khẩn trương tìm những người sống sót trong này đi!''

``Mau thành lập tổ điều tra vụ việc này, nhanh lên!''

Một giọng trầm khàn khàn bên cạnh cất lên, kèm với đó là những tiếng ho khan: ``Saikyou\dots\ *khụ khụ* Em\dots\ vẫn ổn chứ?''

Tôi giật mình, quay sang, thấy Nii-chan cũng đang khom người, tóc rối tung vì bụi, vài vệt máu mảnh loang trên trán.

``\dots Em\dots\ không sao\dots'' Tôi nói như không thật sự tin vào chính câu trả lời ấy. Bằng một cách kỳ diệu nào đó, hai anh em chúng tôi vẫn còn sống.

``Hết\dots *khụ* động đất chưa vậy\dots\ Nii-chan?'' tôi hỏi, cảm thấy họng càng rát hơn vì bụi.

``Có vẻ là\dots\ vãn rồi đấy. Mau rời khỏi đây nào.''

Chúng tôi cùng chậm rãi bò ra khỏi gầm bàn. Chân tôi run rẩy khi vừa duỗi ra khỏi tư thế co quắp. Khung cảnh trước mắt khiến tôi muốn nôn.

Nơi Hanabi-san và Inari-san vừa đứng\dots\ giờ chỉ còn lại một đống đổ nát cao ngang ngực. Những mảng bê tông, gạch vỡ và thanh sắt cong vênh xếp chồng hỗn loạn lên nhau. Cái khoảng trống từng là bàn giáo viên cũng sụp hoàn toàn, bụi vẫn còn bay lơ lửng như khói.


Tôi bật gọi: ``\textbf{Hanabi-san! Inari-san! Saki-san! Touka-san! Suzune-san! Có ai nghe thấy không!?}''

Không một tiếng trả lời.

Nii-chan cũng cất tiếng, giọng anh mạnh mẽ hơn nhưng ẩn chứa sự căng thẳng khó giấu: ``\textbf{Mọi người! Nếu còn nghe được thì lên tiếng đi!}''

Chỉ có âm vọng khô khốc của chính anh vang lại từ những bức tường bị xô lệch và đổ nát.

Tim tôi siết lại. Không thể nào\dots\ Không thể nào tất cả họ\dots

Tôi lao tới một khoảng đổ nát, cố dùng tay gạt bớt gạch vụn, vừa bới vừa lạc giọng gọi: ``\textbf{Có ai còn ở đó không!? Mau trả lời bọn mình đi!}''

Nhưng Nii-chan kéo tôi lại, bàn tay anh siết chặt vado cánh tay tôi đến mức đau đớn. ``Saikyou! Đừng liều lĩnh! Đống này có thể sập tiếp bất cứ lúc nào!''

``Nhưng họ\dots!''

``Anh biết,'' anh khẽ ngắt lời, ánh mắt nghiêm lại. ``Nhưng trước tiên chúng ta phải tìm lối ra. Nếu ở đây lâu hơn, cả hai cũng sẽ bị chôn vùi.''

Tôi im bặt, không thể cãi lại. Lý trí hiểu anh nói đúng, nhưng trái tim vẫn đau buốt. Chúng tôi cùng nhau khom người, lần mò trong đống hoang tàn của lớp 1-C, từng bước lách qua những mảnh gỗ gãy, những bàn ghế gãy đôi như xương vỡ. Mỗi tiếng động tôi tạo ra đều vang lên quá lớn trong cái khoảng không chết chóc này.

Phải mất vài phút, chúng tôi mới thấy khung cửa lớp, nửa trên đã sập, chỉ còn một khoảng trống thấp sát đất. Cả hai phải cúi rạp người mới chui ra được.

Và rồi\dots\ cả hai chúng tôi đông cứng.

Bên ngoài, ở hành lang lớp học.

Sạch sẽ. Nguyên vẹn.

Không một vết nứt, không một hạt bụi nào. Ánh nắng chiều vẫn rọi qua ô cửa kính như chưa từng có chuyện gì xảy ra. Không khí mát và trong đến vô lý so với cái hỗn loạn ngột ngạt phía sau lưng.

\img{e1-3g}

Tôi quay phắt lại. Phía sau cánh cửa kia của lớp 1-C vẫn là địa ngục bụi đá. Nhưng phía trước chúng tôi giờ đây là hành lang tươi sáng, yên bình như thể thời gian chưa từng dao động.

``\dots Chuyện\dots\ quái gì vừa xảy ra thế này?'' tôi thốt ra, giọng run như sắp gãy.

Nii-chan đứng sát bên, lặng vài giây rồi mới nói, thấp giọng đến mức gần như thì thầm: ``Giống\dots\ giống hệt như trong giấc mơ\dots\ Giống như nơi đầu tiên mà anh em mình nhìn thấy ở đó\dots''

Câu nói ấy khiến tôi lạnh buốt từ gáy xuống lưng. Tôi cảm tưởng rằng\dots\ mọi thứ đang diễn ra theo một cách rất quen thuộc.

Chúng tôi bước đi dọc hành lang, và cảm giác bất an càng lúc càng rõ rệt. Mỗi bước chân vang lên đều như bị nuốt chửng bởi khoảng không, không một tiếng vọng trở lại. Khung cảnh này\dots\ tôi đã thấy nó rồi, trong giấc mơ đêm trước ngày nhập học. Từng ô cửa, từng tấm bảng xanh ở bên ngoài, đến cả những chi tiết nhỏ ở trên trần lát ô đều khớp nguyên xi.

Chỉ có một điều khác biệt: Bên trong mỗi lớp học.

Có lớp, qua ô cửa kính, là một đống đất đá cao ngất, như thể cả căn phòng vừa bị nuốt chửng bởi sạt lở. Có lớp lại bị che khuất hoàn toàn bởi một màu đen đặc quánh, như mực, không phản chiếu ánh sáng. Nhìn vào lâu khiến tôi hoa mắt, cảm giác như mình sẽ rơi vào đó nếu cứ tiếp tục dán mắt thêm vài giây nữa.

Nii-chan và tôi chỉ trao đổi nhau những cái nhìn ngắn ngủi, không ai muốn nói thành lời điều đang cùng nghĩ tới: Chúng tôi đang đi lại y nguyên con đường của giấc mơ.

Cuối cùng, hành lang kết thúc.

Ở đó, đứng quay lưng về phía chúng tôi là một cô gái tóc nâu dài đến lưng, mặc đồng phục đen với chiếc khăn quàng xanh. Mái tóc cô khẽ lay động trong luồng gió mỏng từ ô cửa cuối hành lang.

Tôi nhận ra ngay. Nii-chan cũng nhận ra.

Misora Shino.

Người xuất hiện thoáng qua trong lớp 1-C trước khi động đất xảy ra.

Người đã đứng trước mặt chúng tôi trong giấc mơ, giữa hành lang vô tận này, nhưng ở một thời kỳ khác.

``Shino-san\dots?'' giọng Nii-chan vang lên trước khi tôi kịp mở miệng.

Cô gái khẽ quay lại. Ánh mắt xanh của cô lướt qua hai chúng tôi, hơi ngạc nhiên nhưng lại nhanh chóng trở về vẻ điềm tĩnh, như thể đây không phải là lần đầu tiên gặp mặt.

\img{e1-3h}

``Vậy\dots'' cô khẽ nói, giọng đều đều nhưng đủ để từng từ rạch vào không khí, ``\dots hai người có vẻ như là các học sinh mới chuyển đến à.''

Cụm từ ``học sinh mới chuyển'' khiến cả tôi và Nii-chan thoáng rùng mình. Trong giấc mơ, chúng tôi đã đọc những dòng nhật ký cảnh báo, những lời đồn, và những báo cáo liên quan tới cụm từ này. Inari cũng đã nhắc tới một lời truyền miệng tương tự ngoài đời. Mọi mảnh ghép bắt đầu xếp lại thành một hình thù đáng sợ.

Shino nhìn chúng tôi thêm một lúc, rồi khóe môi cong lên. Không phải nụ cười chào hỏi. Không phải nụ cười vui vẻ. Mà là thứ gì đó vừa nham hiểm, vừa như đang thử thách.

Cơn ớn lạnh chạy dọc sống lưng tôi từ đêm hôm ấy tới giờ\dots\ vẫn không hề biến mất.

\end{document}